\chapter{Nucleofiele substitutie}\label{ch:b1:nucleophilic-substitution}

Neem een fles \cip{S}-2-broombutaan, één enkel \omterm{def:b1:stereochemistry-in-depth:enantiomers}{enantiomeer} dat
gepolariseerd licht draait. Verwarmd met natriumhydroxide in een polair
\omterm{def:b1:intermolecular-forces-solvents:solvent-classes}{aprotisch oplosmiddel} geeft het butaan-2-ol dat ook één enkel \omterm{def:b1:stereochemistry-in-depth:enantiomers}{enantiomeer}
is --- maar het andere, \cip{R}: de reactie heeft het molecuul
binnenstebuiten gekeerd, zoals een paraplu in de wind. Laat men hetzelfde
bromide daarentegen in warm water staan, dan geeft het butaan-2-ol dat
gedeeltelijk racemisch is. Dezelfde totale verandering, een OH-groep die
een Br-atoom vervangt, volgt twee verschillende wegen. Dit hoofdstuk
stelt de twee mechanismen vast uit hun \omterm{def:b1:rate-laws:order}{snelheidswetten}, verklaart hun
stereochemie, en geeft de regels die tussen hen beslissen --- het eerste
volledige gebruik van de werktuigen van de hoofdstukken over kinetiek en
elektronische effecten.

\begin{recall}
\omterm{def:b1:rate-laws:order}{Snelheidswetten}, \omterm{def:b1:elementary-steps:elementary-step}{elementaire stappen}, de \omterm{def:b1:elementary-steps:rds}{snelheidsbepalende stap} en de
\omterm{def:b1:elementary-steps:qssa}{stationaire-toestandsbenadering} staan in
\cref{ch:b1:rate-laws,ch:b1:elementary-steps}; \omterm{def:b1:stereochemistry-in-depth:isomers}{configuraties},
weergaven volgens Cram en \omterm{def:b1:stereochemistry-in-depth:optical-activity}{optische activiteit} in
\cref{ch:b1:stereochemistry-in-depth}; \omterm{def:b1:electronic-effects:nucleophile}{nucleofielen}, \omterm{def:b1:electronic-effects:nucleophile}{elektrofielen},
\omterm{def:b1:electronic-effects:nucleophile}{vertrekkende groepen} en \omterm{def:b1:electronic-effects:intermediates}{carbokationen} in \cref{ch:b1:electronic-effects};
protische en \omterm{def:b1:intermolecular-forces-solvents:solvent-classes}{aprotische oplosmiddelen} in
\cref{ch:b1:intermolecular-forces-solvents}.
\end{recall}

\section{De substitutie en haar partners}

\begin{definition}[Nucleofiele substitutie, substraat]\label{def:b1:nucleophilic-substitution:substitution}
Een \emph{nucleofiele substitutie}\index{nucleofiele substitutie} is de
vervanging, op een verzadigd koolstofatoom, van een \omterm{def:b1:electronic-effects:nucleophile}{vertrekkende groep} Y
door een \omterm{def:b1:electronic-effects:nucleophile}{nucleofiel} Nu: $\mathrm{Nu^-} + \mathrm{R{-}Y} \to \mathrm{R{-}Nu} +
\mathrm{Y^-}$. De verbinding R--Y is het
\emph{substraat}\index{substraat}; het koolstofatoom dat Y draagt, is
\omterm{def:b1:electronic-effects:nucleophile}{elektrofiel} omdat de C--Y-binding naar Y toe gepolariseerd is.
\end{definition}

Halogeenalkanen zijn de typische \omterm{def:b1:nucleophilic-substitution:substitution}{substraten}, met \ce{Cl-}, \ce{Br-} of
\ce{I-} als \omterm{def:b1:electronic-effects:nucleophile}{vertrekkende groepen}. Hydroxide geeft alcoholen; alkoxiden
geven ethers; cyanide geeft nitrillen (één koolstofatoom meer); ammoniak
en aminen geven aminen; jodide wisselt uit met de andere halogeniden.
Water en alcoholen, gebruikt als oplosmiddel, kunnen ook als \omterm{def:b1:electronic-effects:nucleophile}{nucleofiel}
optreden: de substitutie is dan een \omterm{def:b1:nucleophilic-substitution:sn1}{solvolyse}.

\section{Het SN2-mechanisme}

Wanneer broommethaan met hydroxide in water reageert, verdubbelt een
verdubbeling van elk van beide concentraties de snelheid:
$v = k[\ce{CH3Br}][\ce{OH-}]$, een reactie van orde één in elke beginstof.

\begin{definition}[SN2-mechanisme, waldeninversie]\label{def:b1:nucleophilic-substitution:sn2}
Het \emph{SN2-mechanisme}\index{SN2-mechanisme} (substitutie,
\omterm{def:b1:electronic-effects:nucleophile}{nucleofiel}, bimoleculair) is één enkele \omterm{def:b1:elementary-steps:elementary-step}{elementaire stap} waarin het
\omterm{def:b1:electronic-effects:nucleophile}{nucleofiel} het koolstofatoom aanvalt aan de kant tegenover de vertrekkende
groep, terwijl de \omterm{def:b1:electronic-effects:nucleophile}{vertrekkende groep} vertrekt. De drie andere
substituenten van het koolstofatoom klappen om, zoals een paraplu die
binnenstebuiten waait: de
\emph{waldeninversie}\index{waldeninversie}.
\end{definition}

\begin{omfigure}
\setchemfig{atom sep=2.1em}
\schemestart
\chemfig{H@{nu}O^{-}}
\arrow{0}[,0.35]
\chemfig{@{c}C(-[:150]H_3C)(<[:210]H)(<:[:250]C_2H_5)-[@{cb}]@{br}Br}
\arrow{->}
\chemfig{[:0]HO-C(-[:30]CH_3)(<[:-30]H)(<:[:-70]C_2H_5)}
\+
\chemfig{Br^{-}}
\schemestop
\chemmove{\draw[->, thick, shorten <=2pt, shorten >=3pt](nu) .. controls +(15:5mm) and +(175:5mm) .. (c);
          \draw[->, thick, shorten <=2pt](cb) .. controls +(90:6mm) and +(90:6mm) .. (br);}

\medskip
\begin{tikzpicture}[font=\small]
  \node[draw, dashed, rounded corners, inner sep=6pt] at (0,0)
    {$\left[\vcenter{\hbox{\ce{HO}$\cdots$\chemfig{C(-[2]CH_3)(-[:-30]H)(-[:210]C_2H_5)}$\cdots$\ce{Br}}}\right]^{-}$};
  \node[below] at (0,-1.1) {overgangstoestand: vijf groepen rond C, drie in een vlak};
\end{tikzpicture}

{\small De SN2-reactie van hydroxide met \cip{S}-2-broombutaan. Het
\omterm{def:b1:electronic-effects:nucleophile}{nucleofiel} valt de zijde tegenover het broom aan (\omterm{def:b1:electronic-effects:curly-arrow}{gebogen pijlen}); in de
\omterm{def:b1:elementary-steps:energy-profile}{overgangstoestand} wordt de O--C-binding gevormd terwijl de C--Br-binding
breekt; het product is \cip{R}-butaan-2-ol: de drie andere groepen zijn
naar de andere kant omgeklapt.}
\end{omfigure}

\begin{proposition}[Snelheidswet van SN2]\label{prop:b1:nucleophilic-substitution:sn2-law}
Voor een SN2-reactie is $v = k[\mathrm{RY}][\mathrm{Nu}]$.
\end{proposition}

\begin{proof}
Het mechanisme is één enkele bimoleculaire \omterm{def:b1:elementary-steps:elementary-step}{elementaire stap}; volgens de
massawerkingswet voor \omterm{def:b1:elementary-steps:elementary-step}{elementaire stappen} (\cref{ch:b1:elementary-steps})
is de snelheid ervan evenredig met het product van de concentraties van
de twee partners.
\end{proof}

\begin{proposition}[Stereochemie van SN2]\label{prop:b1:nucleophilic-substitution:sn2-inversion}
Een SN2-reactie op een stereogeen koolstofatoom geeft één enkel
\omterm{def:b1:stereochemistry-in-depth:isomers}{stereo-isomeer}, met een omgekeerde schikking van de drie ongewijzigde
groepen: een reactie van dit soort is stereospecifiek.
\end{proposition}

\begin{proof}
Het \omterm{def:b1:electronic-effects:nucleophile}{nucleofiel} kan alleen de lege kant van het koolstofatoom bereiken,
tegenover de \omterm{def:b1:electronic-effects:nucleophile}{vertrekkende groep}; de drie andere groepen gaan in de
\omterm{def:b1:elementary-steps:energy-profile}{overgangstoestand} door een vlak en eindigen aan de andere kant. De nieuwe
binding wijst waar de oude niet wees: elk molecuul \omterm{def:b1:nucleophilic-substitution:substitution}{substraat} geeft de
gespiegelde schikking van zijn groepen. De descriptor verandert meestal
van \cip{S} in \cip{R} of omgekeerd, maar niet altijd, omdat hij afhangt
van de rangen van het \omterm{def:b1:electronic-effects:nucleophile}{nucleofiel} en de \omterm{def:b1:electronic-effects:nucleophile}{vertrekkende groep}.
\end{proof}

\begin{definition}[Stereospecifieke reactie]\label{def:b1:nucleophilic-substitution:stereospecific}
Een reactie is een
\emph{stereospecifieke reactie}\index{stereospecifieke reactie}
wanneer stereo-isomere \omterm{def:b1:nucleophilic-substitution:substitution}{substraten} \omterm{def:b1:stereochemistry-in-depth:isomers}{stereo-isomeer} verschillende producten
geven, elk \omterm{def:b1:nucleophilic-substitution:substitution}{substraat} zijn eigen product.
\end{definition}

\section{Het SN1-mechanisme}

2-Broom-2-methylpropaan reageert met water met een snelheid die niet
afhangt van de concentratie van het \omterm{def:b1:electronic-effects:nucleophile}{nucleofiel}: $v = k[(\ce{CH3})_3\ce{CBr}]$.
Een stap waarbij het \omterm{def:b1:electronic-effects:nucleophile}{nucleofiel} niet betrokken is, moet de snelheid
bepalen.

\begin{definition}[SN1-mechanisme, racemisatie, solvolyse]\label{def:b1:nucleophilic-substitution:sn1}
Het \emph{SN1-mechanisme}\index{SN1-mechanisme} (substitutie, \omterm{def:b1:electronic-effects:nucleophile}{nucleofiel},
unimoleculair) heeft twee stappen: de trage \omterm{def:b1:electronic-effects:cleavage}{heterolyse} van de
C--Y-binding, die een \omterm{def:b1:electronic-effects:intermediates}{carbokation} geeft, en daarna de snelle additie van
het \omterm{def:b1:electronic-effects:nucleophile}{nucleofiel} aan het \omterm{def:b1:electronic-effects:intermediates}{carbokation}. Een reactie die één enkel \omterm{def:b1:stereochemistry-in-depth:enantiomers}{enantiomeer}
omzet in een mengsel van beide, is een
\emph{racemisatie}\index{racemisatie}; een substitutie waarin het
oplosmiddel het \omterm{def:b1:electronic-effects:nucleophile}{nucleofiel} is, is een \emph{solvolyse}\index{solvolyse}.
\end{definition}

\begin{omfigure}
\setchemfig{atom sep=2.1em}
\schemestart
\chemfig{C(-[2]CH_3)(-[4]H_3C)(-[6]CH_3)-[@{cb2}]@{br2}Br}
\arrow{->[traag]}
\chemfig{C^{+}(-[2]CH_3)(-[:210]H_3C)(-[:-30]CH_3)}
\+
\chemfig{Br^{-}}
\schemestop

\smallskip
\schemestart
\chemfig{C^{+}(-[2]CH_3)(-[:210]H_3C)(-[:-30]CH_3)}
\arrow{->[snel][\ce{H2O}]}
\chemfig{C(-[2]CH_3)(-[4]H_3C)(-[6]CH_3)-OH_2^{+}}
\arrow{->[\ce{-H+}]}
\chemfig{C(-[2]CH_3)(-[4]H_3C)(-[6]CH_3)-OH}
\schemestop
\chemmove{\draw[->, thick, shorten <=2pt](cb2) .. controls +(90:6mm) and +(90:6mm) .. (br2);}

\medskip
\begin{tikzpicture}[font=\small]
  \fill[omProp!15, draw=omProp] (0,0.2) ellipse (0.2 and 0.85);
  \fill[omProp!15, draw=omProp] (0,-0.2) ellipse (0.2 and 0.85);
  \draw[thick] (0,0) -- (-1.1,-0.2) node[left] {\ce{CH3}};
  \draw[thick] (0,0) -- (0.95,-0.45) node[right] {\ce{C2H5}};
  \draw[thick] (0,0) -- (0.4,0.5) node[above right] {H};
  \node[fill=white, inner sep=1pt] at (0,0) {\ce{C+}};
  \draw[->, thick, omDef] (-0.9,1.6) node[left] {\ce{H2O}} -- (-0.15,0.95);
  \draw[->, thick, omDef] (0.9,-1.6) node[right] {\ce{H2O}} -- (0.15,-0.95);
  \node[align=center] at (4.6,0) {aanval op beide zijden:\\ beide configuraties,\\ racemisch als niets anders\\ de twee zijden onderscheidt};
\end{tikzpicture}

{\small SN1-solvolyse van 2-broom-2-methylpropaan in water: trage
ionisatie tot het tertiaire \omterm{def:b1:electronic-effects:intermediates}{carbokation}, daarna invangen door water en
verlies van een proton. Onder: een \omterm{def:b1:stereochemistry-in-depth:stereocentre}{chiraal} secundair \omterm{def:b1:electronic-effects:intermediates}{carbokation} (uit
2-broombutaan) is vlak, en water addeert met gelijke kans aan beide
zijden.}
\end{omfigure}

\begin{proposition}[Snelheidswet van SN1]\label{prop:b1:nucleophilic-substitution:sn1-law}
Voor het \omterm{def:b1:nucleophilic-substitution:sn1}{SN1-mechanisme} is, als het \omterm{def:b1:electronic-effects:intermediates}{carbokation} veel sneller ingevangen
wordt dan het naar het \omterm{def:b1:nucleophilic-substitution:substitution}{substraat} terugkeert, $v = k_1[\mathrm{RY}]$,
onafhankelijk van het \omterm{def:b1:electronic-effects:nucleophile}{nucleofiel}.
\end{proposition}

\begin{proof}
Stappen: $\mathrm{RY} \rightleftharpoons \mathrm{R^+} + \mathrm{Y^-}$ ($k_1$,
$k_{-1}$), $\mathrm{R^+} + \mathrm{Nu} \to \mathrm{RNu}$ ($k_2$). Het
\omterm{def:b1:electronic-effects:intermediates}{carbokation} is een \omterm{def:b1:electronic-effects:intermediates}{reactief intermediair}: volgens de
\omterm{def:b1:elementary-steps:qssa}{stationaire-toestandsbenadering} (\cref{ch:b1:elementary-steps}) is
$k_1[\mathrm{RY}] =
k_{-1}[\mathrm{R^+}][\mathrm{Y^-}] + k_2[\mathrm{R^+}][\mathrm{Nu}]$, dus
$v = k_2[\mathrm{R^+}][\mathrm{Nu}] = \dfrac{k_1k_2[\mathrm{RY}][\mathrm{Nu}]}
{k_{-1}[\mathrm{Y^-}] + k_2[\mathrm{Nu}]}$. Als $k_2[\mathrm{Nu}] \gg
k_{-1}[\mathrm{Y^-}]$, herleidt dit zich tot $k_1[\mathrm{RY}]$: de eerste
stap is snelheidsbepalend.
\end{proof}

\begin{proposition}[Stereochemie van SN1]\label{prop:b1:nucleophilic-substitution:sn1-stereo}
Een SN1-reactie op een stereogeen koolstofatoom geeft beide
\omterm{def:b1:stereochemistry-in-depth:isomers}{configuraties}: de reactie is niet stereospecifiek, en leidt tot
\omterm{def:b1:nucleophilic-substitution:sn1}{racemisatie} wanneer niets de twee zijden van het \omterm{def:b1:electronic-effects:intermediates}{carbokation}
onderscheidt.
\end{proposition}

\begin{proof}
Het \omterm{def:b1:electronic-effects:intermediates}{carbokation} is vlak (\cref{ch:b1:electronic-effects}): zijn twee
zijden zijn elkaars spiegelbeeld. Is de \omterm{def:b1:electronic-effects:nucleophile}{vertrekkende groep} vóór het
invangen ver weg, dan addeert het \omterm{def:b1:electronic-effects:nucleophile}{nucleofiel} met dezelfde snelheid aan
beide zijden, wat de twee \omterm{def:b1:stereochemistry-in-depth:enantiomers}{enantiomeren} in gelijke hoeveelheden geeft. In
de praktijk schermt het vertrekkende ion nog een tijdje één zijde af, en
wordt vaak een lichte overmaat aan inversie waargenomen.
\end{proof}

\begin{omfigure}
\begin{tikzpicture}
\begin{axis}[
  width=0.8\linewidth, height=5.0cm, xmin=0, xmax=10, ymin=-30, ymax=110,
  axis x line=bottom, axis y line=left, xlabel={reactiecoördinaat},
  ylabel={energie}, xtick=\empty, ytick=\empty,
  label style={font=\small}, legend style={font=\scriptsize, draw=none, at={(0.98,0.98)}, anchor=north east}]
\addplot[omDef, very thick] table[x=x, y=E] {figdata/out/bachelor-1-energy-profiles-sn2.dat};
\addplot[omProp, very thick, dashed] table[x=x, y=E] {figdata/out/bachelor-1-energy-profiles-sn1.dat};
\legend{SN2: één stap, SN1: twee stappen}
\node[font=\scriptsize, anchor=north] at (axis cs:5,68) {carbokation};
\node[font=\scriptsize, anchor=north] at (axis cs:0.6,-3) {RY};
\node[font=\scriptsize, anchor=south] at (axis cs:9.6,-18) {RNu};
\end{axis}
\end{tikzpicture}

{\small Schematische \omterm{def:b1:elementary-steps:energy-profile}{energieprofielen}. SN2: één enkele \omterm{def:b1:elementary-steps:energy-profile}{overgangstoestand},
vijf groepen rond koolstof. SN1: een eerste, hogere barrière (ionisatie,
de \omterm{def:b1:elementary-steps:rds}{snelheidsbepalende stap}), het \omterm{def:b1:electronic-effects:intermediates}{carbokation} in een ondiep dal, en daarna
een kleine barrière voor het invangen ervan.}
\end{omfigure}

\section{Wat het mechanisme bepaalt}

\begin{proposition}[Substraat]\label{prop:b1:nucleophilic-substitution:substrate}
Methyl- en primaire \omterm{def:b1:nucleophilic-substitution:substitution}{substraten} reageren volgens SN2; tertiaire \omterm{def:b1:nucleophilic-substitution:substitution}{substraten}
volgens SN1; secundaire \omterm{def:b1:nucleophilic-substitution:substitution}{substraten} volgens een van beide, afhankelijk van
de andere factoren. Allylische en benzylische \omterm{def:b1:nucleophilic-substitution:substitution}{substraten} reageren snel
volgens beide.
\end{proposition}

\begin{proof}
SN2 heeft ruimte achter het koolstofatoom nodig: elke alkylgroep in plaats
van een waterstofatoom maakt de \omterm{def:b1:elementary-steps:energy-profile}{overgangstoestand}, die vijf groepen
draagt, voller en vertraagt de reactie; drie groepen blokkeren haar. SN1
heeft een stabiel \omterm{def:b1:electronic-effects:intermediates}{carbokation} nodig: de stabiliteitsvolgorde
(\cref{prop:b1:electronic-effects:carbocation-order}) maakt de ionisatie
snel voor tertiaire, mogelijk voor secundaire, en te traag voor primaire
en methylcarbokationen. Allylische en benzylische \omterm{def:b1:electronic-effects:intermediates}{carbokationen} worden
door resonantie gestabiliseerd, en hetzelfde $\pi$-systeem stabiliseert
de \omterm{def:b1:elementary-steps:energy-profile}{overgangstoestand} van SN2.
\end{proof}

\begin{proposition}[Vertrekkende groep]\label{prop:b1:nucleophilic-substitution:leaving-group}
Beide mechanismen zijn sneller naarmate de \omterm{def:b1:electronic-effects:nucleophile}{vertrekkende groep} beter is,
dat wil zeggen naarmate de base die ze wordt zwakker is:
$\ce{I-} > \ce{Br-} > \ce{Cl-} \gg \ce{F-}$; \ce{OH-} en \ce{RO-}
vertrekken niet.
\end{proposition}

\begin{proof}
In beide wordt de C--Y-binding verbroken in of vóór de
\omterm{def:b1:elementary-steps:rds}{snelheidsbepalende stap}, en draagt de \omterm{def:b1:elementary-steps:energy-profile}{overgangstoestand} een deel van de
negatieve lading op Y. Een groep die een negatieve lading goed draagt ---
de geconjugeerde base van een \omterm{def:b1:predominant-reaction:strong-acid}{sterk zuur} --- verlaagt die
\omterm{def:b1:elementary-steps:energy-profile}{overgangstoestand}. \ce{HI}, \ce{HBr} en \ce{HCl} zijn \omterm{def:b1:predominant-reaction:strong-acid}{sterke zuren},
\ce{HF} een \omterm{def:b1:predominant-reaction:strong-acid}{zwak zuur} ($\mathrm{p}K_a$ 3.2), water en alcoholen zeer
zwakke: \ce{OH-} en \ce{RO-} zijn \omterm{def:b1:predominant-reaction:strong-acid}{sterke basen} en slechte vertrekkende
groepen.
\end{proof}

\begin{proposition}[Oplosmiddel]\label{prop:b1:nucleophilic-substitution:solvent}
Polaire \omterm{def:b1:intermolecular-forces-solvents:solvent-classes}{protische oplosmiddelen} (water, alcoholen) begunstigen SN1;
polaire \omterm{def:b1:intermolecular-forces-solvents:solvent-classes}{aprotische oplosmiddelen} (propanon, dimethylsulfoxide,
dimethylformamide) begunstigen SN2 met anionische \omterm{def:b1:electronic-effects:nucleophile}{nucleofielen}.
\end{proposition}

\begin{proof}
SN1 maakt twee ionen uit een neutraal \omterm{def:b1:nucleophilic-substitution:substitution}{substraat}: een polair \omterm{def:b1:intermolecular-forces-solvents:solvent-classes}{protisch
oplosmiddel} stabiliseert beide, het kation via zijn \omterm{def:b1:lewis-resonance-vsepr:lewis}{vrije paren} en het
anion via \omterm{def:b1:intermolecular-forces-solvents:hbond}{waterstofbruggen}, en verlaagt de ionisatiebarrière. Bij SN2
wordt de lading van een anionisch \omterm{def:b1:electronic-effects:nucleophile}{nucleofiel} in de \omterm{def:b1:elementary-steps:energy-profile}{overgangstoestand}
gespreid; een \omterm{def:b1:intermolecular-forces-solvents:solvent-classes}{protisch oplosmiddel}, dat het kleine anion sterk via
\omterm{def:b1:intermolecular-forces-solvents:hbond}{waterstofbruggen} bindt, vertraagt de reactie, terwijl een \omterm{def:b1:intermolecular-forces-solvents:solvent-classes}{aprotisch
oplosmiddel} het anion vrij en reactief laat
(\cref{ch:b1:intermolecular-forces-solvents}).
\end{proof}

\begin{method}[Het mechanisme kiezen]\label{met:b1:nucleophilic-substitution:choose-mechanism}
\begin{enumerate}
  \item \omterm{def:b1:nucleophilic-substitution:substitution}{Substraat}: methyl of primair $\to$ SN2; tertiair $\to$ SN1;
        secundair $\to$ ga verder.
  \item \omterm{def:b1:electronic-effects:nucleophile}{Nucleofiel}: sterk en anionisch ($\ce{OH-}$, $\ce{RO-}$, $\ce{CN-}$,
        $\ce{I-}$) $\to$ SN2; zwak en neutraal (water, alcohol) $\to$ SN1.
  \item Oplosmiddel: aprotisch $\to$ SN2; protisch $\to$ SN1.
  \item Voorspel de stereochemie (inversie of \omterm{def:b1:nucleophilic-substitution:sn1}{racemisatie}) en controleer
        dat een \omterm{def:b1:predominant-reaction:strong-acid}{sterke base} niet veeleer een eliminatie veroorzaakt
        (\cref{ch:b1:elimination}).
\end{enumerate}
\end{method}

\begin{omfigure}
\begin{tabular}{lll}
\hline
 & SN2 & SN1 \\
\hline
stappen & één & twee, via een \omterm{def:b1:electronic-effects:intermediates}{carbokation} \\
\omterm{def:b1:rate-laws:order}{snelheidswet} & $k[\mathrm{RY}][\mathrm{Nu}]$ & $k[\mathrm{RY}]$ \\
\omterm{def:b1:nucleophilic-substitution:substitution}{substraat} & methyl $>$ primair $>$ secundair & tertiair $>$ secundair \\
\omterm{def:b1:electronic-effects:nucleophile}{nucleofiel} & sterk, anionisch & zwak, neutraal (vaak het oplosmiddel) \\
oplosmiddel & polair aprotisch & polair protisch \\
stereochemie & inversie (stereospecifiek) & \omterm{def:b1:nucleophilic-substitution:sn1}{racemisatie} (meestal) \\
\hline
\end{tabular}

\smallskip
{\small De twee mechanismen naast elkaar.}
\end{omfigure}

\section{Oefeningen}

\begin{exercise}[$\star$]\label{exo:b1:nucleophilic-substitution:1}
Schrijf de producten op van: broomethaan met natriumhydroxide;
joodmethaan met natriumethoxide; 1-chloorpropaan met kaliumcyanide;
broommethaan met ammoniak (eerste stap). Wijs \omterm{def:b1:nucleophilic-substitution:substitution}{substraat}, \omterm{def:b1:electronic-effects:nucleophile}{nucleofiel} en
\omterm{def:b1:electronic-effects:nucleophile}{vertrekkende groep} aan.
\end{exercise}

\begin{exercise}[$\star$]\label{exo:b1:nucleophilic-substitution:2}
Wat gebeurt er voor een reactie $\mathrm{RY} + \mathrm{Nu}$ met
$v = k[\mathrm{RY}][\mathrm{Nu}]$ met de \omterm{def:b1:rate-laws:initial-rate}{beginsnelheid} als
$[\mathrm{Nu}]$ verdrievoudigd wordt? Als beide concentraties gehalveerd
worden? Dezelfde vragen als $v = k[\mathrm{RY}]$.
\end{exercise}

\begin{exercise}[$\star$]\label{exo:b1:nucleophilic-substitution:3}
Voorspel het mechanisme (SN1 of SN2): 1-broombutaan met natriumjodide in
propanon; 2-broom-2-methylpropaan in water; broommethaan met hydroxide;
2-broombutaan in ethanol zonder toegevoegde base.
\end{exercise}

\begin{exercise}[$\star$]\label{exo:b1:nucleophilic-substitution:4}
Teken het product van de SN2-reactie van cyanide met
\cip{R}-2-broombutaan in de \omterm{def:b1:stereochemistry-in-depth:representations}{weergave volgens Cram} en geef zijn
descriptor.
\end{exercise}

\begin{exercise}[$\star\star$]\label{exo:b1:nucleophilic-substitution:5}
Schrijf het volledige \omterm{def:b1:nucleophilic-substitution:sn1}{SN1-mechanisme} van de hydrolyse van
2-broom-2-methylpropaan op, met \omterm{def:b1:electronic-effects:curly-arrow}{gebogen pijlen}, inclusief het verlies van
het proton. Waarom is de \omterm{def:b1:rate-laws:order}{snelheidswet} van de eerste orde, hoewel er drie
deeltjes aan deelnemen?
\end{exercise}

\begin{exercise}[$\star\star$]\label{exo:b1:nucleophilic-substitution:6}
Leg uit waarom 1-broom-2,2-dimethylpropaan, een primair bromide, zeer
traag volgens SN2 reageert en niet volgens SN1.
\end{exercise}

\begin{exercise}[$\star\star$]\label{exo:b1:nucleophilic-substitution:7}
\cip{S}-2-joodoctaan verliest zijn \omterm{def:b1:stereochemistry-in-depth:optical-activity}{optische activiteit} wanneer het in
propanon met natriumjodide bewaard wordt, hoewel er geen nieuwe
verbinding ontstaat. Verklaar, en zeg waarom de snelheid waarmee de
activiteit verdwijnt twee keer de substitutiesnelheid is.
\end{exercise}

\begin{exercise}[$\star\star$]\label{exo:b1:nucleophilic-substitution:8}
Rangschik naar snelheid van SN2 met jodide: 2-broom-2-methylpropaan,
broommethaan, 2-broompropaan, 1-broompropaan. Rangschik dezelfde
verbindingen naar snelheid van SN1-solvolyse.
\end{exercise}

\begin{exercise}[$\star\star$]\label{exo:b1:nucleophilic-substitution:9}
Een optisch zuiver secundair bromide ondergaat \omterm{def:b1:nucleophilic-substitution:sn1}{solvolyse}; het product
draait het licht over \qty{12}{\percent} van de waarde die voor zuivere
inversie verwacht wordt. Geef de verhoudingen van de twee \omterm{def:b1:stereochemistry-in-depth:enantiomers}{enantiomeren}
van het product en interpreteer.
\end{exercise}

\begin{exercise}[$\star\star\star$]\label{exo:b1:nucleophilic-substitution:10}
Leid de algemene \omterm{def:b1:rate-laws:order}{snelheidswet} van SN1 af met de
\omterm{def:b1:elementary-steps:qssa}{stationaire-toestandsbenadering}, en toon aan dat het toevoegen van het
ion van de \omterm{def:b1:electronic-effects:nucleophile}{vertrekkende groep}, \ce{Y-}, aan het begin de reactie
vertraagt (het \omterm{def:b1:precipitation:common-ion}{effect van het gemeenschappelijke ion} in de kinetiek).
Onder welke voorwaarde is dit effect verwaarloosbaar?
\end{exercise}

\begin{exercise}[$\star\star\star$]\label{exo:b1:nucleophilic-substitution:11}
Twee \omterm{def:b1:rate-laws:half-life}{halveringstijden} voor de reactie van broommethaan met hydroxide: toon
met $[\ce{OH-}]_0 = [\ce{CH3Br}]_0 = C_0$ aan dat $1/[\ce{CH3Br}] = 1/C_0 +
kt$, en geef $t_{1/2}$. Toon met $[\ce{OH-}]_0 = 50\,C_0$ aan dat de
afname exponentieel is, en geef $t_{1/2}$. Gegevens van de oefening: $k =
\qty{2.0e-4}{L/(mol.s)}$, $C_0 = \qty{0.010}{mol/L}$.
\end{exercise}

\begin{exercise}[$\star\star\star$]\label{exo:b1:nucleophilic-substitution:12}
Leg uit waarom een alcohol, \ce{ROH}, niet reageert met natriumbromide,
maar wel met geconcentreerd broomwaterstofzuur tot \ce{RBr}. Schrijf het
mechanisme op voor een tertiaire en voor een primaire alcohol.
\end{exercise}

\section{Probleem: solvolyse van tert-butylchloride}

\begin{problem}[{Weekendprobleem --- orde en snelheidsconstante uit
geleidbaarheidsgegevens, de halveringstijd, het mechanisme en zijn
stereochemie, en de activeringsenergie van de solvolyse}]
\label{pb:b1:nucleophilic-substitution:1}
2-Chloor-2-methylpropaan (tert-butylchloride) wordt, bij een kleine
concentratie, opgelost in een mengsel van water en propanon. Zijn
\omterm{def:b1:nucleophilic-substitution:sn1}{solvolyse} \ce{(CH3)3CCl + H2O -> (CH3)3COH + H+ + Cl-} vormt ionen, en het
geleidingsvermogen $\kappa$ van de oplossing, dat evenredig is met de
hoeveelheid gevormde ionen, wordt geregistreerd. Na zeer lange tijd
bereikt het bij beide temperaturen $\kappa_\infty = \qty{2.000}{mS/cm}$.
Metingen (\unit{mS/cm}):

\begin{center}
\begin{tabular}{lcccccccc}
\hline
$t$ (min) & 0 & 5 & 10 & 15 & 20 & 30 & 45 & 60 \\
\hline
$\kappa$, \qty{25.0}{\celsius} & 0 & 0.172 & 0.329 & 0.473 & 0.605 & 0.835 & 1.110 & 1.321 \\
$\kappa$, \qty{35.0}{\celsius} & 0 & 0.507 & 0.885 & 1.168 & 1.379 & 1.654 & 1.856 & 1.940 \\
\hline
\end{tabular}
\end{center}
% data generated in figdata/bachelor-1/solvolysis.py (story data)

\medskip
\noindent\textbf{Deel I --- De orde.}
\begin{enumerate}
  \item Waarom is het geleidingsvermogen evenredig met de hoeveelheid
        \omterm{def:b1:nucleophilic-substitution:substitution}{substraat} die gereageerd heeft?
  \item Druk $[\mathrm{RCl}]/[\mathrm{RCl}]_0$ uit als functie van $\kappa$
        en $\kappa_\infty$.
  \item Schrijf de geïntegreerde wet van een eerste-ordereactie op en de
        grootheid die tegen $t$ uitgezet moet worden om haar te toetsen.
  \item Bereken $\ln(\kappa_\infty - \kappa)$ bij \qty{25.0}{\celsius} voor
        elk tijdstip.
  \item Is de grafiek een rechte? Trek een besluit over de orde.
  \item Water is in grote overmaat. Wat verbergt dat in de orde?
\end{enumerate}

\medskip
\noindent\textbf{Deel II --- \omterm{def:b1:rate-laws:order}{Snelheidsconstante} en \omterm{def:b1:rate-laws:half-life}{halveringstijd}.}
\begin{enumerate}[resume]
  \item Bepaal de \omterm{def:b1:rate-laws:order}{snelheidsconstante} bij \qty{25.0}{\celsius}, in
        \unit{s^{-1}}.
  \item Bereken de \omterm{def:b1:rate-laws:half-life}{halveringstijd}.
  \item Op welk tijdstip is \qty{90}{\percent} van het \omterm{def:b1:nucleophilic-substitution:substitution}{substraat}
        verbruikt?
  \item Bepaal de \omterm{def:b1:rate-laws:order}{snelheidsconstante} bij \qty{35.0}{\celsius}.
  \item Met welke factor is de snelheid voor tien graden toegenomen?
  \item Controleer één waarde van de reeks bij \qty{35}{\celsius} aan de
        wet.
\end{enumerate}

\medskip
\noindent\textbf{Deel III --- Het mechanisme.}
\begin{enumerate}[resume]
  \item Welk mechanisme suggereren het \omterm{def:b1:nucleophilic-substitution:substitution}{substraat} en het oplosmiddel?
  \item Schrijf het op met \omterm{def:b1:electronic-effects:curly-arrow}{gebogen pijlen}.
  \item Toon aan dat het in overeenstemming is met de gevonden orde.
  \item Natriumhydroxide toevoegen tot \qty{0.01}{mol/L} verandert de
        \omterm{def:b1:rate-laws:rate}{verdwijningssnelheid} van het chloride nauwelijks. Waarom?
  \item Hetzelfde experiment met het \omterm{def:b1:stereochemistry-in-depth:stereocentre}{chirale} tertiaire chloride
        3-chloor-3-methylhexaan, optisch zuiver, geeft een alcohol met
        bijna geen \omterm{def:b1:stereochemistry-in-depth:optical-activity}{optische activiteit}. Verklaar.
  \item Teken een \omterm{def:b1:elementary-steps:energy-profile}{energieprofiel} van de reactie en duid de
        \omterm{def:b1:elementary-steps:rds}{snelheidsbepalende stap} aan.
  \item Waarom zou 1-chloorbutaan onder dezelfde omstandigheden nauwelijks
        reageren?
\end{enumerate}

\medskip
\noindent\textbf{Deel IV --- \omterm{def:b1:rate-laws:activation-energy}{Activeringsenergie}.}
\begin{enumerate}[resume]
  \item Schrijf de \omterm{thm:b1:rate-laws:arrhenius}{wet van Arrhenius} op.
  \item Druk $E_a$ uit met de twee \omterm{def:b1:rate-laws:order}{snelheidsconstanten}.
  \item Bereken $E_a$ ($R = \qty{8.314}{J/(K.mol)}$).
  \item Wat meet deze energie in het mechanisme?
  \item Voorspel de \omterm{def:b1:rate-laws:order}{snelheidsconstante} bij \qty{45.0}{\celsius}.
  \item Geef de \omterm{def:b1:rate-laws:activation-energy}{activeringsenergie} van de \omterm{def:b1:nucleophilic-substitution:sn1}{solvolyse}, op twee beduidende
        cijfers.
\end{enumerate}
\end{problem}
